% Responsável: TODO(grupo)
% Enunciado, item 6.3: como os dados foram divididos, como as features foram
% extraídas, como o modelo foi treinado/ajustado.
\chapter{Metodologia}
\label{cap:metodologia}

\section{Segmentação em janelas}
\label{sec:janelas}

TODO(grupo): o tamanho da janela é \emph{derivado} da rotação do eixo e da taxa
de amostragem, de modo a cobrir várias revoluções. Apresentar a conta, não só o
número escolhido.

\section{Protocolo de divisão treino/teste}
\label{sec:split}

TODO(grupo): divisão por condição de carga --- treino com \SIlist{0;1;2}{\hp},
teste com \SI{3}{\hp}. Explicar por que a divisão aleatória em nível de janela é
inválida aqui: janelas vizinhas da mesma gravação são quase idênticas e o modelo
passa a memorizar a gravação, não a falha.

\section{Extração de features}
\label{sec:features}

TODO(grupo): domínio do tempo, domínio da frequência, envelope.

\section{Modelos}
\label{sec:modelos}

\subsection{Detecção}
TODO(grupo): treinada \emph{somente} com dados normais.

\subsection{Classificação}
TODO(grupo).

\section{Reprodutibilidade}
\label{sec:reprodutibilidade}

TODO(grupo): seeds fixas e declaradas, versões das bibliotecas, ambiente de
execução (devcontainer e Colab).
